\chapter{Практика 2. Помехоустойчивое кодирование}

\section{Задание}
Реализовать помехоустойчивое кодирование с использованием кода Хэмминга.

\section{Теоретическое описание}
Код Хэмминга — это линейный блочный код, позволяющий обнаруживать и исправлять одиночные ошибки. Для кода Хэмминга с $k$ информационными битами количество контрольных битов $m$ определяется из условия:
\[
2^m \geq k + m + 1
\]
Общая длина кодового слова: $n = k + m$. Контрольные биты располагаются на позициях, являющихся степенями двойки (1, 2, 4, 8, ...). Каждый контрольный бит вычисляется как XOR (сложение по модулю 2) информационных битов, номера позиций которых имеют единицу в соответствующем разряде двоичного представления.

\section{Реализация}

\begin{lstlisting}[language=Python, caption=Кодек Хэмминга]
class HammingCoder:
    def __init__(self, k_bits):
        self.K = k_bits
        self.m = 0
        while (2 ** self.m) < (self.K + self.m + 1):
            self.m += 1
        self.N = self.K + self.m

    def encode(self, bits):
        if not bits:
            return bits
        remainder = len(bits) % self.K
        if remainder:
            bits = bits + '0' * (self.K - remainder)
        encoded = []
        for i in range(0, len(bits), self.K):
            data = [int(b) for b in bits[i:i + self.K]]
            codeword = [0] * self.N
            data_idx = 0
            for pos in range(1, self.N + 1):
                if not (pos & (pos - 1) == 0):
                    codeword[pos - 1] = data[data_idx]
                    data_idx += 1
            for j in range(self.m):
                p_pos = 2 ** j
                parity = 0
                for pos in range(1, self.N + 1):
                    if pos & p_pos:
                        parity ^= codeword[pos - 1]
                codeword[p_pos - 1] = parity
            encoded.extend([str(b) for b in codeword])
        return ''.join(encoded)

    def decode(self, bits):
        if not bits or len(bits) % self.N != 0:
            return bits if not bits else None
        decoded = []
        for i in range(0, len(bits), self.N):
            r = [int(b) for b in bits[i:i + self.N]]
            syndrome = 0
            for j in range(self.m):
                p_pos = 2 ** j
                parity = 0
                for pos in range(1, self.N + 1):
                    if pos & p_pos:
                        parity ^= r[pos - 1]
                if parity:
                    syndrome += p_pos
            if 0 < syndrome <= self.N:
                r[syndrome - 1] ^= 1
            for pos in range(1, self.N + 1):
                if not (pos & (pos - 1) == 0):
                    decoded.append(str(r[pos - 1]))
        return ''.join(decoded)
\end{lstlisting}

\section{Результат}
При использовании кода Хэмминга с 11 информационными битами:
\begin{itemize}
    \item Количество контрольных битов: $m = 4$
    \item Длина кодового слова: $n = 15$
    \item Скорость кодирования: $R = 11/15 \approx 0.733$
\end{itemize}

Исходное сообщение из 408 бит после кодирования занимает 570 бит.

\section{Вывод}
Код Хэмминга успешно реализован. Он позволяет исправлять одиночные ошибки, что повышает помехоустойчивость передачи данных.